\subsection{Matrix Inverses}
\hrulefill

\paragraph*{Multiplicative Inverse}
If $A$ is a square matrix and $AB = BA = I$, then $B$ is called the multiplicative inverse of $A$ and is denoted by $A^{-1} = B$.

\paragraph*{Note} Let $A \in \mathbb{R}^{m \times n} \ni m \neq n$ and $B \in \mathbb{R}^{n \times m}$, it is possible for 
$A \cdot B = I_m$ but $B \cdot A \neq I_n$. In this case, $A$ is called a left inverse of $B$ and $B$ is called a right inverse of $A$.

\paragraph*{Properties}
Let $A\vec{x} = \vec{b}$, then
\begin{itemize}
    \item $A^{-1}A\vec{x} = A^{-1}\vec{b} \implies \vec{x} = A^{-1}\vec{b}$
    \item $I\vec{x} = A^{-1}\vec{b}$
    \item $\vec{x} = A^{-1}\vec{b}$
\end{itemize}

\paragraph*{Theorem} 
If $A$ is invertible, then $A\vec{x} = \vec{0}$ has 1 unique solution.

\paragraph*{Example} Show that $B = \begin{bmatrix}
    -1 & 1 \\
    1 & 0
\end{bmatrix}$ is the inverse of $A = \begin{bmatrix}
    0 & 1 \\
    1 & 1
\end{bmatrix}$

\[AB = \begin{bmatrix}
    0 & 1 \\
    1 & 1
\end{bmatrix} \begin{bmatrix}
    -1 & 1 \\
    1 & 0
\end{bmatrix} = \begin{bmatrix}
    1 & 0 \\
    0 & 1
\end{bmatrix} = I\]
\[BA = \begin{bmatrix}
    -1 & 1 \\
    1 & 0
\end{bmatrix} \begin{bmatrix}
    0 & 1 \\
    1 & 1
\end{bmatrix} = \begin{bmatrix}
    1 & 0 \\
    0 & 1
\end{bmatrix} = I\]

\paragraph*{Example} 
Show the $A = \begin{bmatrix}
    1 & 2 \\
    0 & 0
\end{bmatrix}$ has no inverse.
\[AB = \begin{bmatrix}
    1 & 2 \\
    0 & 0
\end{bmatrix} \begin{bmatrix}
    a & b \\
    c & d
\end{bmatrix} = \begin{bmatrix}
    a + 2c & b + 2d \\
    0 & 0
\end{bmatrix} \neq I\]

\subsubsection*{Determinants of 2x2 Matricies}
Let $A = \begin{bmatrix}
a & b \\
c & d
\end{bmatrix}$, then $\text{det}(A) = ad - bc$. If $\text{det}(A) \neq 0$, then $A$ is invertible and $A^{-1} = \frac{1}{\text{det}(A)} \begin{bmatrix}
d & -b \\
-c & a
\end{bmatrix}$

\paragraph*{Theorem}
$A\vec{x} = \vec{b}$ and $\exists A^{-1}$, then $\vec{x} = A^{-1}\vec{b}$

\paragraph*{Example} % 2.4.1

\subsubsection*{Matrix Inversion Algorithm}
If $A_{n\times n}$, then there is a sequence of elementary row operations that carry $A \rightarrow I_n$.
The same sequence of elementary row operations carries $I_n \rightarrow A^{-1}$.
\begin{enumerate}
    \item Form the augmented matrix $[A | I_n]$.
    \item Use elementary row operations to reduce the left half of the augmented matrix to $I_n$.
    \item The right half of the augmented matrix is now $A^{-1}$.
    \item If the left half of the augmented matrix cannot be reduced to $I_n$, then $A$ is not invertible.
\end{enumerate}

\paragraph*{Example} % 2.4.2

\subsubsection*{Properties of Inverses}
Let $A$ be invertible, then
\begin{itemize}
    \item $AB = AC \implies B = C$
    \item $(A^T)^{-1} = (A^{-1})^T$(
    \item $(A^{-1})^{-1} = A$
    \item $(AB)^{-1} = B^{-1}A^{-1}$
    \item $AB(AB)^{-1} = AB(B^{-1}A^{-1}) = I$
    \item $[A_1\ A_2\ \dots\ A_k]^{-1} = [A_k^{-1}\ \dots\ A_2^{-1}\ A_1^{-1}]$
    \item $(A^k)^{-1} = (A^{-1})^k \forall k \in \mathbb{R}, k > 0$
    \item $(aA)^{-1} = \frac{1}{a}A^{-1} \forall a \in \mathbb{R}, a \neq 0$
\end{itemize}

\paragraph*{Example} % 2.4.3

\paragraph*{Inverse Theorem}
The following conditions are equivilent for an $n \times n$ matrix $A$:
\begin{itemize}
    \item $A$ is invertible
    \item $RREF(A) = I_n$
    \item rank($A$) = $n$
    \item The homogeneous system $A\vec{x} = \vec{0}$ has only the trivial solution $\vec{x} = \vec{0}$
    \item The system $A\vec{x} = \vec{b}$ has a unique solution for every $\vec{b} \in \mathbb{R}^n$
    \item $det(A) \neq 0$
    \item The linear transformation $T(\vec{x}) = A\vec{x}$ is one-to-one and onto.
\end{itemize}

\paragraph*{Inverses of Block Matricies}
Let $P = \begin{bmatrix}
    A_{m \times m} & X \\
    0 & B_{n \times n}
\end{bmatrix}$ and $Q = \begin{bmatrix}
    A_{m \times m} & 0 \\
    Y & B_{n \times n}
\end{bmatrix}$, then $P^{-1} = \begin{bmatrix}
    A^{-1} & -A^{-1}XB^{-1} \\
    0 & B^{-1}
\end{bmatrix}$ and $Q^{-1} = \begin{bmatrix}
    A^{-1} & 0 \\
    -B^{-1}YA^{-1} & B^{-1}
\end{bmatrix}$

\paragraph*{Theorem}
A matrix transformation $T_A$ is invertible if and only if the matrix A is invertable. $T_A$ being invertible means T is onto and one to one.
\[T_A(\vec{x}) = A\vec{x} \implies T_A^{-1}(\vec{y}) = A^{-1}\vec{y}\]